% !TEX root = ../main.tex
%% Chapitre 1 — Raisonnement et vocabulaire ensembliste
%% Programme officiel MPSI/MP2I (BO spécial n°1 du 11/02/2021), premier semestre,
%% section « Raisonnement et vocabulaire ensembliste » :
%%   a) rudiments de logique ; b) ensembles ; c) applications et relations.
%% Résultats numériques et identités vérifiés par verifications/chapitre_01.py.

%% --- Mini-diagrammes de Venn (utilisés dans la synthèse) --------------
\newcommand{\otVenn}[2]{%
  \begin{tikzpicture}
    \draw[otMuted,line width=0.5pt,rounded corners=1.5mm] (-1.85,-1.2) rectangle (1.85,1.2);
    \node[anchor=north west,font=\scriptsize,text=otMuted,inner sep=2pt] at (-1.85,1.2) {$E$};
    #1
    \draw[otNavy,line width=0.7pt] (-0.42,0) circle (0.75);
    \draw[otGreenDark,line width=0.7pt] (0.42,0) circle (0.75);
    \node[font=\scriptsize,text=otNavy] at (-1.0,0.62) {$A$};
    \node[font=\scriptsize,text=otGreenDark] at (1.0,0.62) {$B$};
    \node[anchor=north,font=\otLight\fontsize{7.5}{9}\selectfont,text=otInk] at (0,-1.3) {#2};
  \end{tikzpicture}}

\begin{ouverture}[ensembles]{1}{Raisonnement et vocabulaire ensembliste}{Partie 1 · Premier semestre}
  {Logique\enspace·\enspace Raisonnements\enspace·\enspace Récurrences\enspace·\enspace
   Ensembles\enspace·\enspace Applications\enspace·\enspace Relations}
  \objectif{Manipuler propositions, connecteurs et quantificateurs, et écrire la négation d'un énoncé.}
  \objectif{Choisir le bon mode de raisonnement : contraposée, absurde, disjonction de cas, analyse-synthèse.}
  \objectif{Rédiger une récurrence simple, double ou forte.}
  \objectif{Maîtriser les opérations sur les parties d'un ensemble, le produit cartésien et les partitions.}
  \objectif{Étudier une application : images directes et réciproques, injectivité, surjectivité, bijectivité.}
  \objectif{Reconnaître une relation d'équivalence, une relation d'ordre.}
  \prerequis{Ensembles de nombres usuels $\N$, $\Z$, $\Q$, $\R$ et calcul littéral du lycée.}
  \prerequis{Raisonnement par récurrence rencontré en Terminale.}
  \prerequis{Vocabulaire des fonctions : image, antécédent, courbe représentative.}
\end{ouverture}

%% =====================================================================
%%  L'ESSENTIEL
%% =====================================================================
\essentiel

\essSection{1}{Propositions et quantificateurs}

\begin{definition}[Proposition]
Une \textbf{proposition} est un énoncé mathématique qui est soit vrai, soit faux.
À partir de deux propositions $P$ et $Q$, on forme les propositions
$\text{non}(P)$, «~$P$ et $Q$~», «~$P$ ou $Q$~» (le « ou » est inclusif),
$P\Rightarrow Q$ («~$P$ implique $Q$~») et $P\Leftrightarrow Q$
(«~$P$ équivaut à $Q$~»).
\end{definition}

\noindent
\begin{minipage}[c]{0.56\linewidth}
L'implication $P\Rightarrow Q$ est fausse dans un seul cas : lorsque $P$ est
vraie et $Q$ fausse. En particulier, si $P$ est fausse, l'implication
$P\Rightarrow Q$ est vraie, quelle que soit $Q$.\par\vspace{2mm}
L'équivalence $P\Leftrightarrow Q$ signifie «~$P\Rightarrow Q$ \emph{et}
$Q\Rightarrow P$~» : démontrer une équivalence, c'est démontrer
\emph{deux} implications.
\end{minipage}\hfill
\begin{minipage}[c]{0.4\linewidth}
\centering
{\renewcommand{\arraystretch}{1.3}\arrayrulecolor{otRule}%
\begin{tabular}{@{}cc|cc@{}}
\toprule
$P$ & $Q$ & $P\Rightarrow Q$ & $P\Leftrightarrow Q$ \\
\midrule
V & V & V & V \\
V & F & \textcolor{otRed}{F} & F \\
F & V & V & F \\
F & F & V & V \\
\bottomrule
\end{tabular}}\\[1mm]
\legende{Table de vérité de l'implication et de l'équivalence.}
\end{minipage}

\begin{definition}[Réciproque, contraposée, négation d'une implication]
Soit l'implication $P\Rightarrow Q$.
\begin{itemize}[label={\color{otNavy}\textbullet},leftmargin=1.4em,itemsep=0.6mm,topsep=0.8mm]
  \item Sa \textbf{réciproque} est $Q\Rightarrow P$ : elle n'a, en général, pas la même valeur de vérité.
  \item Sa \textbf{contraposée} est $\text{non}(Q)\Rightarrow\text{non}(P)$ :
    elle est \emph{toujours} équivalente à $P\Rightarrow Q$.
  \item Sa \textbf{négation} est «~$P$ et $\text{non}(Q)$~» — ce n'est pas une implication !
\end{itemize}
\end{definition}

\begin{propriete}[Négations usuelles]
$\text{non}(P\text{ et }Q)$ est «~$\text{non}(P)$ ou $\text{non}(Q)$~» ;
$\text{non}(P\text{ ou }Q)$ est «~$\text{non}(P)$ et $\text{non}(Q)$~» ;
$\text{non}(\text{non}(P))$ est $P$.
\end{propriete}

\otIntertitre{Quantificateurs}
Le quantificateur \textbf{universel} $\forall$ se lit « pour tout » ; le
quantificateur \textbf{existentiel} $\exists$ se lit « il existe (au moins
un) ». Le symbole $\exists\,!$ signifie « il existe un unique ».
La négation échange les deux quantificateurs :
\[
  \text{non}\bigl(\forall x\in E,\ P(x)\bigr)\ :\ \exists x\in E,\ \text{non}\bigl(P(x)\bigr)
  \qquad\qquad
  \text{non}\bigl(\exists x\in E,\ P(x)\bigr)\ :\ \forall x\in E,\ \text{non}\bigl(P(x)\bigr)
\]

\begin{attention}[L'ordre des quantificateurs est essentiel]
On ne peut pas échanger $\forall$ et $\exists$ :
«~$\forall x\in\R,\ \exists y\in\R,\ y>x$~» est vraie (pour chaque $x$, on
peut choisir $y=x+1$, qui \emph{dépend} de $x$), alors que
«~$\exists y\in\R,\ \forall x\in\R,\ y>x$~» est fausse (un même $y$ devrait
dépasser tous les réels, y compris $y$ lui-même).
Par ailleurs, les quantificateurs appartiennent au langage mathématique :
on ne les emploie jamais comme abréviations dans une phrase rédigée.
\end{attention}

\essSection{2}{Les modes de raisonnement}

{\renewcommand{\arraystretch}{1.5}\arrayrulecolor{otRule}%
\begin{tabularx}{\linewidth}{@{}>{\otSemi\fontsize{7.5}{10}\selectfont\leavevmode\color{otNavy}}p{30mm}Y@{}}
\toprule
\otCaps[6]{Direct} & Pour montrer $P\Rightarrow Q$ : on suppose $P$ vraie et on en déduit $Q$ par une chaîne de déductions. \\
\otCaps[6]{Contraposition} & Au lieu de $P\Rightarrow Q$, on démontre $\text{non}(Q)\Rightarrow\text{non}(P)$, qui lui est équivalente. \\
\otCaps[6]{Par l'absurde} & Pour montrer qu'une proposition $P$ est vraie, on suppose $\text{non}(P)$ et on aboutit à une contradiction. \\
\otCaps[6]{Disjonction de cas} & On découpe la situation en plusieurs cas qui recouvrent toutes les possibilités, et on conclut dans chacun d'eux. \\
\otCaps[6]{Analyse-synthèse} & Pour déterminer tous les objets vérifiant une propriété : l'\emph{analyse} suppose qu'un objet convient et en déduit sa forme (condition nécessaire) ; la \emph{synthèse} vérifie que les candidats trouvés conviennent (condition suffisante). \\
\bottomrule
\end{tabularx}}

\begin{definition}[Condition nécessaire, condition suffisante]
Si $P\Rightarrow Q$ est vraie, on dit que $Q$ est une \textbf{condition
nécessaire} de $P$ (pour que $P$ soit vraie, il faut que $Q$ le soit), et que
$P$ est une \textbf{condition suffisante} de $Q$ (pour que $Q$ soit vraie, il
suffit que $P$ le soit).
\end{definition}

\begin{demonstration}[modèle : $\sqrt{2}$ est irrationnel (par l'absurde)]
Supposons $\sqrt{2}\in\Q$ : il existe des entiers $p$ et $q$ premiers entre
eux, avec $q\geq1$, tels que $\sqrt{2}=\dfrac{p}{q}$. En élevant au carré,
$p^{2}=2q^{2}$, donc $p^{2}$ est pair, donc $p$ est pair (par contraposée :
si $p$ était impair, $p^{2}$ serait impair). Écrivons $p=2k$ : alors
$4k^{2}=2q^{2}$, soit $q^{2}=2k^{2}$ : $q$ est pair à son tour. Ainsi $p$ et
$q$ sont tous deux pairs, ce qui contredit le fait qu'ils sont premiers entre
eux. Donc $\sqrt{2}\notin\Q$.
\end{demonstration}

\essSection{3}{Les raisonnements par récurrence}

\begin{theoreme}[Récurrence simple]
Soit $P(n)$ une propriété dépendant d'un entier naturel $n$, et $n_0\in\N$. Si :
\begin{itemize}[label={\color{otNavy}\textbullet},leftmargin=1.4em,itemsep=0.6mm,topsep=0.8mm]
  \item \etape{Initialisation} $P(n_0)$ est vraie ;
  \item \etape{Hérédité} pour tout entier $k\geq n_0$, $P(k)\Rightarrow P(k+1)$ ;
\end{itemize}
alors $P(n)$ est vraie pour tout entier $n\geq n_0$.
\end{theoreme}

\begin{theoreme}[Récurrence double, récurrence forte]
Avec les mêmes notations :
\begin{itemize}[label={\color{otNavy}\textbullet},leftmargin=1.4em,itemsep=0.6mm,topsep=0.8mm]
  \item \etape{Double} si $P(n_0)$ et $P(n_0+1)$ sont vraies, et si pour tout
    $k\geq n_0$, $\bigl(P(k)\text{ et }P(k+1)\bigr)\Rightarrow P(k+2)$, alors
    $P(n)$ est vraie pour tout $n\geq n_0$ ;
  \item \etape{Forte} si $P(n_0)$ est vraie, et si pour tout $k\geq n_0$,
    $\bigl(P(n_0)\text{ et }P(n_0+1)\text{ et }\cdots\text{ et }P(k)\bigr)\Rightarrow P(k+1)$,
    alors $P(n)$ est vraie pour tout $n\geq n_0$.
\end{itemize}
Ces principes sont liés à la propriété fondamentale suivante : \emph{toute
partie non vide de $\N$ possède un plus petit élément}.
\end{theoreme}

\begin{attention}
Dans l'hérédité, on fixe \emph{un} entier $k$ et on suppose $P(k)$ vraie
(hypothèse de récurrence) : on ne suppose jamais « $P(n)$ vraie pour tout
$n$ », ce qui serait supposer la conclusion. L'initialisation est
indispensable : la propriété « $4^{n}+1$ est divisible par $3$ » est
héréditaire, puisque $4^{k+1}+1=4(4^{k}+1)-3$, et pourtant elle est fausse
pour tout entier $n$.
\end{attention}

\essSection{4}{Ensembles et parties}

\begin{definition}[Appartenance, inclusion]
$x\in E$ signifie que $x$ est un \textbf{élément} de l'ensemble $E$.
Un ensemble $A$ est \textbf{inclus} dans $E$, noté $A\subset E$, lorsque tout
élément de $A$ appartient à $E$ : on dit que $A$ est une \textbf{partie} (ou
un sous-ensemble) de $E$. L'ensemble des parties de $E$ se note
$\parties{E}$ ; l'ensemble vide se note $\varnothing$.
Deux ensembles sont égaux lorsqu'ils ont les mêmes éléments :
$A=B\iff\bigl(A\subset B\text{ et }B\subset A\bigr)$ (\emph{double inclusion}).
\end{definition}

\begin{exemple}
Pour $E=\ens{1,2,3}$ : $\parties{E}=\bigl\{\varnothing,\ens{1},\ens{2},\ens{3},
\ens{1,2},\ens{1,3},\ens{2,3},\ens{1,2,3}\bigr\}$, qui compte $2^{3}=8$ éléments.
Attention à distinguer $1\in E$, $\ens{1}\subset E$ et $\ens{1}\in\parties{E}$.
\end{exemple}

\begin{definition}[Opérations sur les parties]
Soient $A$ et $B$ deux parties de $E$.
La \textbf{réunion} $A\cup B$, l'\textbf{intersection} $A\cap B$, la
\textbf{différence} $A\setminus B=\ens{x\in A \mid x\notin B}$ et le
\textbf{complémentaire} $\overline{A}=E\setminus A$ (aussi noté $A^{c}$)
sont définis par :
\[
  x\in A\cup B\iff\bigl(x\in A\text{ ou }x\in B\bigr)
  \qquad\qquad
  x\in A\cap B\iff\bigl(x\in A\text{ et }x\in B\bigr).
\]
$A$ et $B$ sont \textbf{disjoints} lorsque $A\cap B=\varnothing$.
\end{definition}

\begin{center}
\otVenn{\fill[otNavySoft] (-0.42,0) circle (0.75);\fill[otNavySoft] (0.42,0) circle (0.75);}{$A\cup B$}\hfill
\otVenn{\begin{scope}\clip (-0.42,0) circle (0.75);\fill[otNavySoft] (0.42,0) circle (0.75);\end{scope}}{$A\cap B$}\hfill
\otVenn{\fill[otNavySoft] (-0.42,0) circle (0.75);\begin{scope}\clip (-0.42,0) circle (0.75);\fill[white] (0.42,0) circle (0.75);\end{scope}}{$A\setminus B$}\hfill
\otVenn{\fill[otNavySoft,rounded corners=1.5mm] (-1.85,-1.2) rectangle (1.85,1.2);\fill[white] (-0.42,0) circle (0.75);}{$\overline{A}$}
\end{center}

\begin{propriete}[Règles de calcul]
Pour toutes parties $A$, $B$, $C$ de $E$ :
\begin{itemize}[label={\color{otNavy}\textbullet},leftmargin=1.4em,itemsep=0.6mm,topsep=0.8mm]
  \item \etape{Distributivité} $A\cap(B\cup C)=(A\cap B)\cup(A\cap C)$ et
    $A\cup(B\cap C)=(A\cup B)\cap(A\cup C)$ ;
  \item \etape{Lois de De Morgan} $\overline{A\cup B}=\overline{A}\cap\overline{B}$
    et $\overline{A\cap B}=\overline{A}\cup\overline{B}$ ;
  \item $\overline{\overline{A}}=A$ ;\quad $A\setminus B=A\cap\overline{B}$ ;\quad
    $A\subset B\iff\overline{B}\subset\overline{A}$.
\end{itemize}
\end{propriete}

\begin{definition}[Produit cartésien]
Le \textbf{produit cartésien} $E\times F$ est l'ensemble des couples $(x,y)$
avec $x\in E$ et $y\in F$. Plus généralement, $E_1\times\cdots\times E_n$ est
l'ensemble des $n$-uplets $(x_1,\dots,x_n)$ avec $x_i\in E_i$ pour tout $i$.
On note $E^{n}=E\times\cdots\times E$ ($n$ facteurs).
\end{definition}

\noindent
\begin{minipage}[c]{0.58\linewidth}
\begin{definition}[Recouvrement disjoint, partition]
Des parties $A_1,\dots,A_p$ de $E$ forment un \textbf{recouvrement disjoint}
de $E$ lorsqu'elles sont deux à deux disjointes et que leur réunion est $E$.
Si de plus aucune n'est vide, on parle de \textbf{partition} de $E$ : tout
élément de $E$ appartient alors à une et une seule des parties $A_i$.
\end{definition}
\end{minipage}\hfill
\begin{minipage}[c]{0.38\linewidth}
\centering
\begin{tikzpicture}
  \draw[otMuted,line width=0.6pt,rounded corners=1.8mm] (0,0) rectangle (4.4,2.5);
  \draw[otNavy!60,line width=0.6pt] (1.5,0) .. controls (1.2,0.9) and (1.8,1.7) .. (1.4,2.5);
  \draw[otNavy!60,line width=0.6pt] (3.0,0) .. controls (3.3,0.8) and (2.7,1.8) .. (3.1,2.5);
  \draw[otNavy!60,line width=0.6pt] (1.33,1.3) .. controls (2.2,1.0) and (2.6,1.6) .. (3.0,1.25);
  \node[font=\scriptsize,text=otInk] at (0.75,1.25) {$A_1$};
  \node[font=\scriptsize,text=otInk] at (2.2,1.9) {$A_2$};
  \node[font=\scriptsize,text=otInk] at (2.2,0.55) {$A_3$};
  \node[font=\scriptsize,text=otInk] at (3.8,1.25) {$A_4$};
  \node[anchor=south west,font=\scriptsize,text=otMuted,inner sep=2pt] at (0,2.55) {$E$};
\end{tikzpicture}\\[0.5mm]
\legende{Une partition de $E$ en quatre parties.}
\end{minipage}

\essSection{5}{Applications}

\begin{definition}[Application, graphe]
Une \textbf{application} $f:E\to F$ associe à \emph{tout} élément $x$ de $E$
un \emph{unique} élément de $F$, noté $f(x)$. $E$ est l'ensemble de départ,
$F$ l'ensemble d'arrivée. Le \textbf{graphe} de $f$ est la partie
$\Gamma_f=\ens{(x,f(x)) \mid x\in E}$ de $E\times F$.
L'ensemble des applications de $E$ dans $F$ se note $\mathscr{F}(E,F)$ ou
$F^{E}$. Une \textbf{famille} $(x_i)_{i\in I}$ d'éléments de $E$ est une
application de l'ensemble d'indices $I$ dans $E$.
\end{definition}

\begin{definition}[Indicatrice, restriction]
La \textbf{fonction indicatrice} d'une partie $A$ de $E$ est l'application
$\ind{A}:E\to\ens{0,1}$ qui vaut $1$ sur $A$ et $0$ ailleurs.
La \textbf{restriction} de $f:E\to F$ à une partie $A$ de $E$ est
l'application $f|_{A}:A\to F$ qui coïncide avec $f$ sur $A$ ; on dit alors
que $f$ est un \textbf{prolongement} de $f|_{A}$ à $E$.
\end{definition}

\begin{definition}[Image directe, image réciproque]
Soient $f:E\to F$, $A\subset E$ et $B\subset F$.
\begin{itemize}[label={\color{otNavy}\textbullet},leftmargin=1.4em,itemsep=0.6mm,topsep=0.8mm]
  \item L'\textbf{image directe} de $A$ est $f(A)=\ens{f(x) \mid x\in A}\subset F$ :
    l'ensemble des images des éléments de $A$.
  \item L'\textbf{image réciproque} de $B$ est
    $f^{-1}(B)=\ens{x\in E \mid f(x)\in B}\subset E$ :
    l'ensemble des antécédents des éléments de $B$.
\end{itemize}
\end{definition}

\begin{center}
\begin{tikzpicture}[line cap=round]
  % Image directe
  \draw[otNavy,line width=0.7pt] (0,0) ellipse (1.05 and 1.45);
  \draw[otGreenDark,line width=0.7pt] (4.6,0) ellipse (1.05 and 1.45);
  \node[font=\scriptsize,text=otNavy] at (0,1.75) {$E$};
  \node[font=\scriptsize,text=otGreenDark] at (4.6,1.75) {$F$};
  \draw[otMuted,line width=0.6pt,-{Stealth[length=2mm]}] (1.35,1.25) .. controls (2.3,1.7) .. (3.25,1.25);
  \node[font=\small,text=otInk] at (2.3,1.8) {$f$};
  \fill[otNavySoft] (0,-0.25) ellipse (0.62 and 0.8);
  \draw[otNavy,line width=0.6pt] (0,-0.25) ellipse (0.62 and 0.8);
  \node[font=\scriptsize,text=otNavy] at (0,-0.25) {$A$};
  \fill[otGreenTint] (4.6,-0.25) ellipse (0.55 and 0.7);
  \draw[otGreenDark,line width=0.6pt] (4.6,-0.25) ellipse (0.55 and 0.7);
  \node[font=\scriptsize,text=otGreenDark] at (4.6,-0.25) {$f(A)$};
  \draw[otGreen,line width=0.7pt,-{Stealth[length=1.8mm]}] (0.68,-0.1) -- (3.98,-0.1);
  \node[anchor=north,align=center,font=\otLight\fontsize{7.5}{9}\selectfont,text=otInk] at (2.3,-1.7)
    {Image directe : $f(A)$ regroupe\\les images des éléments de $A$.};
  % Image réciproque
  \begin{scope}[shift={(8.6,0)}]
  \draw[otNavy,line width=0.7pt] (0,0) ellipse (1.05 and 1.45);
  \draw[otGreenDark,line width=0.7pt] (4.6,0) ellipse (1.05 and 1.45);
  \node[font=\scriptsize,text=otNavy] at (0,1.75) {$E$};
  \node[font=\scriptsize,text=otGreenDark] at (4.6,1.75) {$F$};
  \draw[otMuted,line width=0.6pt,-{Stealth[length=2mm]}] (1.35,1.25) .. controls (2.3,1.7) .. (3.25,1.25);
  \node[font=\small,text=otInk] at (2.3,1.8) {$f$};
  \fill[otNavySoft] (0,-0.25) ellipse (0.68 and 0.85);
  \draw[otNavy,line width=0.6pt] (0,-0.25) ellipse (0.68 and 0.85);
  \node[font=\scriptsize,text=otNavy] at (0,-0.25) {$f^{-1}(B)$};
  \fill[otGreenTint] (4.6,-0.25) ellipse (0.55 and 0.7);
  \draw[otGreenDark,line width=0.6pt] (4.6,-0.25) ellipse (0.55 and 0.7);
  \node[font=\scriptsize,text=otGreenDark] at (4.6,-0.25) {$B$};
  \draw[otGreen,line width=0.7pt,{Stealth[length=1.8mm]}-] (0.74,-0.1) -- (3.98,-0.1);
  \node[anchor=north,align=center,font=\otLight\fontsize{7.5}{9}\selectfont,text=otInk] at (2.3,-1.7)
    {Image réciproque : $f^{-1}(B)$ regroupe\\les antécédents de $B$.};
  \end{scope}
\end{tikzpicture}
\end{center}

\begin{attention}
La notation $f^{-1}(B)$ a un sens pour \emph{toute} application $f$ : elle ne
suppose pas que $f$ est bijective, et ne désigne pas une application
réciproque. Ainsi, pour $f:x\mapsto x^{2}$ de $\R$ dans $\R$,
$f^{-1}(\ens{4})=\ens{-2,2}$ alors que $f$ n'est pas bijective.
\end{attention}

\begin{definition}[Composition]
Si $f:E\to F$ et $g:F\to G$, la \textbf{composée} $g\circ f:E\to G$ est
définie par $(g\circ f)(x)=g\bigl(f(x)\bigr)$. La composition est associative :
$h\circ(g\circ f)=(h\circ g)\circ f$.
\end{definition}

\begin{definition}[Injection, surjection, bijection]
Soit $f:E\to F$.
\begin{itemize}[label={\color{otNavy}\textbullet},leftmargin=1.4em,itemsep=0.6mm,topsep=0.8mm]
  \item $f$ est \textbf{injective} si tout élément de $F$ a \emph{au plus} un antécédent :
    $\forall (x,x')\in E^{2},\ f(x)=f(x')\Rightarrow x=x'$.
  \item $f$ est \textbf{surjective} si tout élément de $F$ a \emph{au moins} un antécédent :
    $\forall y\in F,\ \exists x\in E,\ y=f(x)$, c'est-à-dire $f(E)=F$.
  \item $f$ est \textbf{bijective} si elle est injective et surjective : tout
    élément de $F$ a \emph{exactement} un antécédent.
\end{itemize}
\end{definition}

\begin{center}
\begin{tikzpicture}[line cap=round,
  pt/.style={circle,fill=otNavy,inner sep=0pt,minimum size=3.2pt},
  ptv/.style={circle,fill=otGreenDark,inner sep=0pt,minimum size=3.2pt},
  fl/.style={otMuted,line width=0.55pt,-{Stealth[length=1.6mm]}}]
  \foreach \dx/\titre in {0/{injective, non surjective},5.3/{surjective, non injective},10.6/{bijective}}{%
    \draw[otNavy,line width=0.6pt] (\dx,0) ellipse (0.55 and 1.15);
    \draw[otGreenDark,line width=0.6pt] (\dx+2.1,0) ellipse (0.55 and 1.15);
    \node[anchor=north,font=\otLight\fontsize{7.5}{9}\selectfont,text=otInk] at (\dx+1.05,-1.35) {\titre};}
  % injective non surjective : 3 -> 4
  \foreach \y in {0.7,0,-0.7}{\node[pt] at (0,\y) {};}
  \foreach \y in {0.85,0.3,-0.3,-0.85}{\node[ptv] at (2.1,\y) {};}
  \draw[fl] (0.1,0.7) -- (2.0,0.85);
  \draw[fl] (0.1,0) -- (2.0,0.3);
  \draw[fl] (0.1,-0.7) -- (2.0,-0.85);
  % surjective non injective : 4 -> 3
  \foreach \y in {0.85,0.3,-0.3,-0.85}{\node[pt] at (5.3,\y) {};}
  \foreach \y in {0.7,0,-0.7}{\node[ptv] at (7.4,\y) {};}
  \draw[fl] (5.4,0.85) -- (7.3,0.7);
  \draw[fl] (5.4,0.3) -- (7.3,0.7);
  \draw[fl] (5.4,-0.3) -- (7.3,0);
  \draw[fl] (5.4,-0.85) -- (7.3,-0.7);
  % bijective : 3 -> 3
  \foreach \y in {0.7,0,-0.7}{\node[pt] at (10.6,\y) {};}
  \foreach \y in {0.7,0,-0.7}{\node[ptv] at (12.7,\y) {};}
  \draw[fl] (10.7,0.7) -- (12.6,0);
  \draw[fl] (10.7,0) -- (12.6,0.7);
  \draw[fl] (10.7,-0.7) -- (12.6,-0.7);
\end{tikzpicture}
\end{center}

\begin{theoreme}[Bijection réciproque]
Si $f:E\to F$ est bijective, il existe une unique application
$f^{-1}:F\to E$, appelée \textbf{réciproque} de $f$, telle que
$f^{-1}\circ f=\mathrm{Id}_E$ et $f\circ f^{-1}=\mathrm{Id}_F$.
Elle est caractérisée par : $y=f(x)\iff x=f^{-1}(y)$.
$f^{-1}$ est elle-même bijective, et $(f^{-1})^{-1}=f$.
Cette notation est compatible avec l'image réciproque : si $f$ est bijective,
l'image réciproque de $B$ par $f$ est l'image directe de $B$ par $f^{-1}$.
\end{theoreme}

\begin{propriete}[Composées]
La composée de deux injections est injective ; la composée de deux
surjections est surjective ; la composée de deux bijections est bijective,
avec $(g\circ f)^{-1}=f^{-1}\circ g^{-1}$ (attention à l'ordre).
\end{propriete}

\begin{savoirfaire}[Étudier l'injectivité et la surjectivité]
\begin{itemize}[label={\color{otGreenDark}\textbullet},leftmargin=1.4em,itemsep=0.6mm,topsep=0.5mm]
  \item \etape{Injectivité} partir de $f(x)=f(x')$ et montrer $x=x'$
    (ou : montrer que $x\neq x'\Rightarrow f(x)\neq f(x')$ par contraposée).
  \item \etape{Surjectivité} partir d'un $y\in F$ quelconque et \emph{résoudre}
    l'équation $f(x)=y$ d'inconnue $x\in E$.
  \item \etape{Bijectivité} résoudre $f(x)=y$ : si cette équation possède une
    unique solution pour chaque $y\in F$, alors $f$ est bijective, et la
    formule obtenue donne $f^{-1}$.
  \item Pour conclure \emph{négativement}, un contre-exemple explicite suffit.
\end{itemize}
\end{savoirfaire}

\essSection{6}{Relations binaires}

\begin{definition}[Relation binaire]
Une \textbf{relation binaire} $\mathcal{R}$ sur un ensemble $E$ associe à
certains couples $(x,y)\in E^{2}$ la mention « $x$ est en relation avec $y$ »,
notée $x\,\mathcal{R}\,y$. Elle est :
\textbf{réflexive} si $\forall x\in E$, $x\,\mathcal{R}\,x$ ;
\textbf{symétrique} si $x\,\mathcal{R}\,y\Rightarrow y\,\mathcal{R}\,x$ ;
\textbf{antisymétrique} si $\bigl(x\,\mathcal{R}\,y$ et $y\,\mathcal{R}\,x\bigr)\Rightarrow x=y$ ;
\textbf{transitive} si $\bigl(x\,\mathcal{R}\,y$ et $y\,\mathcal{R}\,z\bigr)\Rightarrow x\,\mathcal{R}\,z$.
\end{definition}

\begin{definition}[Relation d'équivalence, classes]
Une relation \textbf{réflexive, symétrique et transitive} est une
\textbf{relation d'équivalence}. La \textbf{classe d'équivalence} de $x$ est
$\mathrm{cl}(x)=\ens{y\in E \mid x\,\mathcal{R}\,y}$.
Les classes d'équivalence forment une \emph{partition} de $E$.
\end{definition}

\begin{exemple}[Congruences]
Soit $c>0$ un réel (ou $c\geq1$ un entier). La relation « $x\equiv y\ [c]$ »,
définie par $x-y\in c\,\Z$, est une relation d'équivalence sur $\R$ (ou sur
$\Z$). Par exemple, $17\equiv 2\ [5]$ et, dans $\R$,
$\frac{9\pi}{4}\equiv\frac{\pi}{4}\ [2\pi]$.
\end{exemple}

\begin{definition}[Relation d'ordre]
Une relation \textbf{réflexive, antisymétrique et transitive} est une
\textbf{relation d'ordre}. L'ordre est \textbf{total} si deux éléments
quelconques sont toujours comparables ($x\,\mathcal{R}\,y$ ou
$y\,\mathcal{R}\,x$) ; sinon, il est \textbf{partiel}.
\end{definition}

\begin{exemple}
$\leq$ est un ordre total sur $\R$. L'inclusion $\subset$ est un ordre
partiel sur $\parties{E}$ dès que $E$ a au moins deux éléments :
$\ens{1}$ et $\ens{2}$ ne sont pas comparables. La divisibilité est un ordre
partiel sur $\N^{*}$ : $2$ et $3$ ne se divisent pas l'un l'autre.
\end{exemple}

\begin{attention}[Pièges fréquents]
\begin{itemize}[label={\color{otOrange}\textbullet},leftmargin=1.4em,itemsep=0.6mm,topsep=0.5mm]
  \item La négation de $P\Rightarrow Q$ n'est pas une implication : c'est « $P$ et $\text{non}(Q)$ ».
  \item Ne pas confondre réciproque et contraposée : seule la contraposée équivaut à l'implication.
  \item «~$\forall x,\ \exists y$~» et «~$\exists y,\ \forall x$~» ne disent pas la même chose :
    dans le premier cas, $y$ peut dépendre de $x$.
  \item Distinguer $\in$ et $\subset$ : $x\in E$, mais $\ens{x}\subset E$ et $\ens{x}\in\parties{E}$.
  \item $f(A\cap A')\subset f(A)\cap f(A')$, et l'inclusion peut être stricte :
    l'égalité n'est garantie que si $f$ est injective.
  \item Pour montrer qu'une application n'est pas injective (ou pas surjective),
    un contre-exemple explicite suffit — le rédiger !
\end{itemize}
\end{attention}

%% =====================================================================
%%  EXERCICES
%% =====================================================================
\exercices
\rubrique{Applications}

\begin{exercice}[1]{Vrai ou faux ?}
Pour chacune des propositions suivantes, dire si elle est vraie ou fausse,
et justifier brièvement.
\begin{questions}
  \item $\forall x\in\R,\ x^{2}\geq0$.
  \item $\exists x\in\R,\ x^{2}=-1$.
  \item $\forall x\in\R,\ \exists y\in\R,\ y>x$.
  \item $\exists y\in\R,\ \forall x\in\R,\ y>x$.
  \item $\forall n\in\N,\ n^{2}+n$ est pair.
  \item $\forall x\in\R,\ \bigl(x^{2}=4\Rightarrow x=2\bigr)$.
\end{questions}
\end{exercice}

\begin{exercice}[1]{Écrire avec des quantificateurs}
Soit $f:\R\to\R$ une fonction et $(u_n)_{n\in\N}$ une suite réelle. Écrire
avec des quantificateurs les énoncés suivants.
\begin{questions}
  \item $f$ est majorée sur $\R$.
  \item $f$ n'est pas majorée sur $\R$ (négation du précédent).
  \item La suite $(u_n)$ est croissante.
  \item $f$ s'annule au moins une fois.
  \item $f$ est paire.
  \item Tout réel admet au moins un antécédent par $f$.
\end{questions}
\end{exercice}

\begin{exercice}[1]{Nier des énoncés quantifiés}
Écrire la négation de chacun des énoncés suivants (sans employer de négation
devant un quantificateur).
\begin{questions}
  \item $\forall x\in\R,\ f(x)>0$.
  \item $\exists x\in\intFF{0}{1},\ f(x)=0$.
  \item $\forall y\in\R,\ \exists x\in\R,\ f(x)=y$.
  \item $\forall\varepsilon>0,\ \exists N\in\N,\ \forall n\geq N,\ |u_n-\ell|\leq\varepsilon$.
  \item $\forall x\in\R,\ \bigl(x\in A\Rightarrow x\in B\bigr)$.
\end{questions}
\end{exercice}

\begin{exercice}[1]{Implication, réciproque, contraposée}
On considère, pour un entier naturel $n$ et un réel $x$, les implications :
\[
  (\mathrm{I}_1)\ :\ \text{« si $n$ est multiple de $4$, alors $n$ est pair »}
  \qquad\qquad
  (\mathrm{I}_2)\ :\ \text{« } x\geq1\Rightarrow x^{2}\geq1\text{ »}.
\]
\begin{questions}
  \item Énoncer la réciproque et la contraposée de $(\mathrm{I}_1)$.
    Parmi $(\mathrm{I}_1)$, sa réciproque et sa contraposée, lesquelles sont vraies ?
  \item Mêmes questions pour $(\mathrm{I}_2)$.
  \item Compléter avec « nécessaire », « suffisante » ou « nécessaire et
    suffisante » : être multiple de $4$ est une condition \ldots{} pour qu'un
    entier soit pair ; être pair est une condition \ldots{} pour qu'un entier
    soit multiple de $4$.
\end{questions}
\end{exercice}

\begin{exercice}[1]{Opérations sur les parties}
\begin{questions}
  \item Dans $E=\ens{1,2,\dots,10}$, on pose $A=\ens{1,2,3,4,5}$ et
    $B=\ens{4,5,6,7}$. Déterminer $A\cup B$, $A\cap B$, $A\setminus B$,
    $B\setminus A$, $\overline{A}$, $\overline{B}$, puis vérifier sur cet
    exemple que $\overline{A\cup B}=\overline{A}\cap\overline{B}$.
  \item Dans $\R$, on pose $I=\intFO{-1}{3}$ et $J=\intOF{1}{5}$.
    Déterminer $I\cap J$, $I\cup J$, $I\setminus J$ et $\R\setminus I$.
\end{questions}
\end{exercice}

\begin{exercice}[1]{Produit cartésien et ensemble des parties}
\begin{questions}
  \item On pose $E=\ens{a,b}$ et $F=\ens{1,2,3}$. Écrire en extension
    l'ensemble $E\times F$. Combien a-t-il d'éléments ?
  \item Écrire en extension $\parties{E}$ pour $E=\ens{1,2,3}$.
  \item Les affirmations suivantes sont-elles vraies ?
    $\varnothing\in\parties{E}$ ;\quad $\varnothing\subset\parties{E}$ ;\quad
    $\ens{1}\in\parties{E}$ ;\quad $1\in\parties{E}$ ;\quad
    $\ens{\ens{1}}\subset\parties{E}$.
\end{questions}
\end{exercice}

\begin{exercice}[1]{Images directes et réciproques : calculs}
On considère l'application $f:\R\to\R$, $x\mapsto x^{2}$.
\begin{questions}
  \item Déterminer les images directes $f\bigl(\intFF{1}{2}\bigr)$,
    $f\bigl(\intFF{-2}{1}\bigr)$ et $f(\R)$.
  \item Déterminer les images réciproques $f^{-1}(\ens{4})$,
    $f^{-1}\bigl(\intFF{1}{4}\bigr)$ et $f^{-1}\bigl(\intFF{-2}{-1}\bigr)$.
  \item Soit $g:\R\to\R$, $x\mapsto2x+1$. Déterminer
    $g\bigl(\intFF{0}{1}\bigr)$ et $g^{-1}\bigl(\intFF{0}{1}\bigr)$.
\end{questions}
\end{exercice}

\begin{exercice}[1]{Injective ? surjective ?}
Pour chacune des applications suivantes, dire si elle est injective, si elle
est surjective, si elle est bijective. Justifier (un contre-exemple suffit
pour une réponse négative).
\begin{tasks}(2)
  \task $f_1:\R\to\R$, $x\mapsto2x+1$
  \task $f_2:\R\to\R$, $x\mapsto x^{2}$
  \task $f_3:\R\to\R_{+}$, $x\mapsto x^{2}$
  \task $f_4:\R_{+}\to\R$, $x\mapsto x^{2}$
  \task $f_5:\N\to\N$, $n\mapsto n+1$
  \task $f_6:\Z\to\Z$, $n\mapsto n+1$
\end{tasks}
\end{exercice}

\rubrique{Entraînement}

\begin{exercice}[2]{Raisonner par contraposée}
\begin{questions}
  \item Soit $n\in\N$. Montrer que si $n^{2}$ est pair, alors $n$ est pair.
  \item Soient $x$ et $y$ deux réels. Montrer que si
    $(x+1)(y-1)=(x-1)(y+1)$, alors $x=y$.
  \item Soit $n\in\N$. Montrer que si $n^{2}-1$ n'est pas divisible par $8$,
    alors $n$ est pair.
\end{questions}
\end{exercice}

\begin{exercice}[2]{Raisonner par l'absurde}
\begin{questions}
  \item En s'inspirant de la démonstration du cours pour $\sqrt{2}$,
    montrer que $\sqrt{3}$ est irrationnel. \emph{On pourra utiliser le fait
    que si $3$ divise $p^{2}$, alors $3$ divise $p$.}
  \item Montrer que la somme d'un rationnel et d'un irrationnel est
    irrationnelle.
  \item Montrer qu'il n'existe pas de plus petit réel strictement positif.
\end{questions}
\end{exercice}

\begin{exercice}[2]{Raisonner par disjonction de cas}
\begin{questions}
  \item Montrer que, pour tout $n\in\N$, $\dfrac{n(n+1)}{2}$ est un entier.
  \item Montrer que, pour tout $n\in\Z$, $3$ divise $n(n^{2}+2)$.
  \item Résoudre dans $\R$ l'équation $|x-1|+|x+2|=5$.
\end{questions}
\end{exercice}

\begin{exercice}[2]{Raisonner par analyse-synthèse}
\begin{questions}
  \item Montrer que toute fonction $f:\R\to\R$ s'écrit de manière unique
    comme la somme d'une fonction paire et d'une fonction impaire.
  \item Déterminer toutes les fonctions $f:\R\to\R$ telles que :
    \[ \forall(x,y)\in\R^{2},\qquad f(x)\,f(y)-f(xy)=x+y. \]
\end{questions}
\end{exercice}

\begin{exercice}[2]{Récurrences simples}
Démontrer par récurrence les propriétés suivantes.
\begin{questions}
  \item Pour tout entier $n\geq1$ :
    $\displaystyle\sum_{k=1}^{n}k^{2}=\frac{n(n+1)(2n+1)}{6}$.
  \item Pour tout entier naturel $n$ : $2^{n}>n$.
  \item Pour tout entier $n\geq4$ : $n!\geq2^{n}$.
  \item Pour tout entier naturel $n$ : $7$ divise $3^{2n+1}+2^{n+2}$.
\end{questions}
\end{exercice}

\begin{exercice}[2]{Récurrence double, récurrence forte}
\begin{questions}
  \item On définit la suite $(u_n)$ par $u_0=2$, $u_1=5$ et, pour tout
    $n\in\N$, $u_{n+2}=5u_{n+1}-6u_n$. Montrer que, pour tout $n\in\N$,
    $u_n=2^{n}+3^{n}$.
  \item Montrer, par récurrence forte, que tout entier $n\geq2$ admet au
    moins un diviseur premier.
  \item On définit la suite $(v_n)$ par $v_1=1$ et, pour tout $n\geq1$,
    $v_{n+1}=v_1+v_2+\dots+v_n$. Calculer $v_2$, $v_3$, $v_4$, puis montrer
    que, pour tout $n\geq2$, $v_n=2^{\,n-2}$.
\end{questions}
\end{exercice}

\begin{exercice}[2]{Calculs ensemblistes}
Soient $A$, $B$ et $C$ trois parties d'un ensemble $E$. Démontrer les
égalités et équivalences suivantes.
\begin{questions}
  \item $A\setminus(A\setminus B)=A\cap B$.
  \item $A\cup(B\cap C)=(A\cup B)\cap(A\cup C)$.
  \item $A\setminus(B\cup C)=(A\setminus B)\cap(A\setminus C)$.
  \item $A\cup B=A\cap B\iff A=B$.
\end{questions}
\end{exercice}

\begin{exercice}[2]{La différence symétrique}
Pour $A$ et $B$ parties d'un ensemble $E$, on pose
$A\,\triangle\,B=(A\setminus B)\cup(B\setminus A)$.
\begin{questions}
  \item Montrer que $A\,\triangle\,B=(A\cup B)\setminus(A\cap B)$.
  \item Déterminer $A\,\triangle\,\varnothing$, $A\,\triangle\,A$ et
    $A\,\triangle\,E$.
  \item Montrer que $A\,\triangle\,B=\varnothing\iff A=B$.
  \item Dans $E=\ens{1,2,\dots,8}$, calculer
    $\ens{1,2,3,4}\,\triangle\,\ens{3,4,5,6}$.
\end{questions}
\end{exercice}

\begin{exercice}[2]{La fonction indicatrice}
Soient $A$ et $B$ deux parties d'un ensemble $E$. On rappelle que
$\ind{A}:E\to\ens{0,1}$ vaut $1$ sur $A$ et $0$ ailleurs.
\begin{questions}
  \item Montrer que $A=B\iff\ind{A}=\ind{B}$, et que
    $A\subset B\iff\ind{A}\leq\ind{B}$.
  \item Montrer que $\ind{A\cap B}=\ind{A}\,\ind{B}$ et que
    $\ind{\overline{A}}=1-\ind{A}$.
  \item En déduire que $\ind{A\cup B}=\ind{A}+\ind{B}-\ind{A}\,\ind{B}$.
  \item Retrouver, à l'aide des questions précédentes, la loi de De Morgan
    $\overline{A\cup B}=\overline{A}\cap\overline{B}$.
\end{questions}
\end{exercice}

\begin{exercice}[2]{Images directes et réciproques : propriétés}
Soient $f:E\to F$ une application, $A$ et $A'$ deux parties de $E$, $B$ et
$B'$ deux parties de $F$.
\begin{questions}
  \item Montrer que si $A\subset A'$, alors $f(A)\subset f(A')$.
  \item Montrer que $f(A\cup A')=f(A)\cup f(A')$.
  \item Montrer que $f(A\cap A')\subset f(A)\cap f(A')$, puis donner un
    exemple où cette inclusion est stricte.
  \item Montrer que $f^{-1}(B\cup B')=f^{-1}(B)\cup f^{-1}(B')$ et que
    $f^{-1}(B\cap B')=f^{-1}(B)\cap f^{-1}(B')$.
  \item Montrer que $A\subset f^{-1}\bigl(f(A)\bigr)$ et que
    $f\bigl(f^{-1}(B)\bigr)\subset B$, et donner un exemple où chacune de ces
    inclusions est stricte.
\end{questions}
\end{exercice}

\rubrique{Approfondissement}

\begin{exercice}[3]{Composition, injectivité et surjectivité}
Soient $f:E\to F$ et $g:F\to G$ deux applications.
\begin{questions}
  \item Montrer que si $f$ et $g$ sont injectives, alors $g\circ f$ est
    injective ; que si $f$ et $g$ sont surjectives, alors $g\circ f$ est
    surjective.
  \item Montrer que si $g\circ f$ est injective, alors $f$ est injective.
  \item Montrer que si $g\circ f$ est surjective, alors $g$ est surjective.
  \item Donner un exemple où $g\circ f$ est bijective, mais où $f$ n'est pas
    surjective et $g$ n'est pas injective.
  \item Soient $f:E\to F$ et $g:F\to E$ telles que $g\circ f=\mathrm{Id}_E$.
    Montrer que $f$ est injective et que $g$ est surjective.
\end{questions}
\end{exercice}

\begin{exercice}[3]{Bijections explicites}
\begin{questions}
  \item On définit $f:\N\to\Z$ par : $f(n)=\dfrac{n}{2}$ si $n$ est pair, et
    $f(n)=-\dfrac{n+1}{2}$ si $n$ est impair. Calculer $f(0)$, $f(1)$,
    $f(2)$, $f(3)$, $f(4)$, puis montrer que $f$ est bijective.
  \item Montrer que $g:\R\setminus\ens{1}\to\R\setminus\ens{2}$,
    $x\mapsto\dfrac{2x+1}{x-1}$, est bijective et expliciter sa réciproque.
  \item Montrer que $h:\R\to\intOO{-1}{1}$,
    $x\mapsto\dfrac{x}{1+|x|}$, est bijective et expliciter sa réciproque.
\end{questions}
\end{exercice}

\begin{exercice}[3]{Équations ensemblistes}
Soient $A$ et $B$ deux parties d'un ensemble $E$. Les inconnues $X$, $Y$ des
équations suivantes sont des parties de $E$.
\begin{questions}
  \item Résoudre le système : $A\cap X=\varnothing$ et $A\cup X=E$.
  \item On suppose $A\subset B$. Montrer que les solutions de l'équation
    $A\cup X=B$ sont exactement les parties $X$ de $E$ telles que
    $B\setminus A\subset X\subset B$. Que se passe-t-il si
    $A\not\subset B$ ?
  \item Montrer que si $A\cap X=A\cap Y$ et $A\cup X=A\cup Y$, alors $X=Y$.
\end{questions}
\end{exercice}

\begin{exercice}[3]{Relations d'équivalence}
\begin{questions}
  \item Sur $\R$, on définit la relation $\mathcal{R}$ par :
    $x\,\mathcal{R}\,y\iff x-y\in\Z$.
    \begin{questions}
      \item Montrer que $\mathcal{R}$ est une relation d'équivalence.
      \item Décrire la classe de $0$, puis celle de $\frac{1}{2}$.
      \item Montrer que chaque classe d'équivalence contient un et un seul
        élément de $\intFO{0}{1}$.
    \end{questions}
  \item Sur $\R$, on définit la relation $\mathcal{S}$ par :
    $x\,\mathcal{S}\,y\iff x^{2}=y^{2}$. Montrer que $\mathcal{S}$ est une
    relation d'équivalence et décrire la classe d'un réel $x$.
\end{questions}
\end{exercice}

\begin{exercice}[3]{La divisibilité, un ordre partiel}
Sur $\N^{*}$, on considère la relation de divisibilité : $a\mid b$ lorsqu'il
existe $k\in\N^{*}$ tel que $b=ka$.
\begin{questions}
  \item Montrer que $\mid$ est une relation d'ordre sur $\N^{*}$.
  \item Cet ordre est-il total ?
  \item Montrer que $1$ est le plus petit élément de $(\N^{*},\mid)$.
    L'ensemble $\N^{*}$ possède-t-il un plus grand élément pour cet ordre ?
  \item On munit $\parties{E}$ de l'inclusion, pour $E$ ayant au moins deux
    éléments. Vérifier que $\subset$ est une relation d'ordre partiel, et
    préciser le plus petit et le plus grand élément de
    $\bigl(\parties{E},\subset\bigr)$.
  \item Dans $(\N^{*},\mid)$, l'ensemble $\ens{4,6}$ admet-il un plus petit
    majorant ? \emph{On décrira l'ensemble des majorants de $\ens{4,6}$.}
\end{questions}
\end{exercice}

\begin{exercice}[3]{Le théorème de Cantor}
L'objectif de ce problème est de montrer qu'il n'existe pas de surjection
d'un ensemble $E$ sur l'ensemble de ses parties $\parties{E}$ — autrement
dit, que $\parties{E}$ est toujours « strictement plus gros » que $E$.
\partieExo{A}{Observations sur un ensemble fini}
On pose $E=\ens{1,2,3}$, et on considère l'application $f:E\to\parties{E}$
définie par $f(1)=\ens{1,2}$, $f(2)=\varnothing$ et $f(3)=\ens{1,3}$.
\begin{questions}
  \item L'application $f$ est-elle injective ? surjective ?
  \item On pose $D=\ens{x\in E \mid x\notin f(x)}$. Déterminer $D$, puis
    vérifier que $D$ n'a pas d'antécédent par $f$.
  \item Combien d'applications de $E$ dans $\parties{E}$ existe-t-il ?
    Combien faudrait-il d'éléments à l'ensemble d'arrivée pour qu'une
    surjection soit possible ?
\end{questions}
\partieExo{B}{Le cas général}
Soit maintenant $E$ un ensemble quelconque et $f:E\to\parties{E}$ une
application. On pose, comme ci-dessus,
$D=\ens{x\in E \mid x\notin f(x)}$.
\begin{questions}[resume]
  \item On raisonne par l'absurde, en supposant $f$ surjective : il existe
    alors $a\in E$ tel que $f(a)=D$. Montrer que l'hypothèse $a\in D$ conduit
    à une contradiction, puis que l'hypothèse $a\notin D$ conduit également à
    une contradiction.
  \item Conclure : énoncer le théorème démontré.
  \item En déduire qu'il n'existe pas de surjection de $\N$ sur
    $\parties{\N}$. \emph{Ce résultat, dû à Georg Cantor (1891), montre qu'il
    existe « plusieurs infinis » : l'infini de $\parties{\N}$ est strictement
    plus grand que celui de $\N$.}
\end{questions}
\end{exercice}

% !TEX root = ../../main.tex
%% Défi du chapitre 1 — Points fixes et théorème de Cantor-Bernstein
%% Résultats vérifiés par verifications/defis_01_05.py.

\rubrique{Défi}

\begin{defi}{Points fixes et théorème de Cantor-Bernstein}{1 h 30}
Soit $E$ un ensemble. Une application $\varphi:\parties{E}\to\parties{E}$
est dite \emph{croissante} si, pour toutes parties $X,Y$ de $E$,
$X\subseteq Y\implies\varphi(X)\subseteq\varphi(Y)$. Le but du problème est
de montrer qu'une telle application possède toujours un point fixe, puis
d'en tirer un théorème célèbre : s'il existe une injection de $E$ dans $F$
et une injection de $F$ dans $E$, alors $E$ et $F$ sont en bijection.

\partieExo{A}{Le théorème du point fixe de Knaster-Tarski}
Soit $\varphi:\parties{E}\to\parties{E}$ croissante. On pose
\[
\mathcal{A}=\ens{X\in\parties{E}\mid X\subseteq\varphi(X)}
\qquad\text{et}\qquad
A=\bigcup_{X\in\mathcal{A}}X .
\]
\begin{questions}
  \item Montrer que $\mathcal{A}$ n'est pas vide.
  \item Soit $X\in\mathcal{A}$. Montrer que $X\subseteq\varphi(A)$ ; en
    déduire que $A\subseteq\varphi(A)$.
  \item Montrer que $\varphi(A)\in\mathcal{A}$, puis que $\varphi(A)=A$.
  \item Montrer que $A$ est le \emph{plus grand} point fixe de $\varphi$ :
    tout $Y$ vérifiant $\varphi(Y)=Y$ est inclus dans $A$.
  \item L'hypothèse de croissance est-elle utile ? Étudier
    $X\mapsto E\setminus X$ pour $E$ non vide.
\end{questions}

\partieExo{B}{Le théorème de Cantor-Bernstein}
Soient $E$, $F$ deux ensembles, $f:E\to F$ et $g:F\to E$ deux applications
\emph{injectives}. Pour $X\subseteq E$, on pose
$\varphi(X)=E\setminus g\bigl(F\setminus f(X)\bigr)$.
\begin{questions}[resume]
  \item Montrer que $\varphi$ est croissante.
  \item Soit $A$ un point fixe de $\varphi$. Montrer que
    $E\setminus A=g\bigl(F\setminus f(A)\bigr)$.
  \item On définit $h:E\to F$ par $h(x)=f(x)$ si $x\in A$, et
    $h(x)=g^{-1}(x)$ si $x\in E\setminus A$. Justifier que $h$ est bien
    définie, puis montrer que $h$ est bijective.
\end{questions}

\medskip
\begin{center}
\begin{tikzpicture}[line cap=round]
  \foreach \xx/\hh/\bb/\nom in {0/A/E\setminus A/E, 5/f(A)/F\setminus f(A)/F}{%
    \begin{scope}[xshift=\xx cm]
      \begin{scope}
        \clip[rounded corners=3mm] (0,0) rectangle (2.5,3.2);
        \fill[otGreen!16] (0,1.55) rectangle (2.5,3.2);
      \end{scope}
      \draw[otNavy,line width=0.8pt,rounded corners=3mm] (0,0) rectangle (2.5,3.2);
      \draw[otNavy,line width=0.5pt,dash pattern=on 2pt off 1.5pt] (0,1.55) -- (2.5,1.55);
      \node[font=\small,text=otGreenDark] at (1.25,2.38) {$\hh$};
      \node[font=\small,text=otNavy] at (1.25,0.78) {$\bb$};
      \node[font=\small,text=otNavy,above] at (1.25,3.22) {$\nom$};
    \end{scope}}
  \draw[otGreenDark,line width=1pt,-{Stealth[length=2.2mm]}]
    (2.6,2.38) to[bend left=14] node[above,font=\small] {$f$} (4.9,2.38);
  \draw[otNavy,line width=1pt,-{Stealth[length=2.2mm]}]
    (4.9,0.78) to[bend left=14] node[below,font=\small] {$g$} (2.6,0.78);
\end{tikzpicture}\\[0.5mm]
\legende{Le découpage de Cantor-Bernstein : $f$ envoie $A$ sur $f(A)$, et
$g$ envoie le reste de $F$ exactement sur le reste de $E$.}
\end{center}

\partieExo{C}{Applications}
\begin{questions}[resume]
  \item On prend $E=\intFF{0}{1}$, $F=\intFO{0}{1}$, $f(x)=x/2$ et
    $g(y)=y$. Déterminer les points fixes de l'application $\varphi$ de la
    partie B, puis expliciter la bijection $h$ obtenue à partir du plus
    grand d'entre eux. Faire un dessin.
  \item En utilisant les injections $n\mapsto(n,0)$ et
    $(p,q)\mapsto2^{p}3^{q}$ (unicité de la décomposition en facteurs
    premiers, chapitre 7), montrer que $\N$ et $\N\times\N$ sont en
    bijection.
  \item Montrer qu'il n'existe aucune injection de $\parties{E}$ dans
    $E$. \emph{On pourra utiliser le théorème de Cantor (exercice 24).}
\end{questions}
\end{defi}

